\documentclass[12pt]{article}
\usepackage[utf8]{inputenc}
\usepackage[margin=0.75in]{geometry}
\usepackage{amsmath}
\usepackage{amssymb}
\usepackage{amsthm}
\usepackage{graphicx} % Required for inserting images
\usepackage{array}
\usepackage{tikz-cd}


\title{MAT1100 Homework 3}
\author{Aaron Avram}
\date{November 1 2026}

\begin{document}
\newtheorem{definition}{Definition}
\newtheorem{theorem}{Theorem}
\newtheorem{claim}{Claim}
\newtheorem{corollary}{Corollary}

\maketitle

\begin{align*}
    \textbf{Question 2}
\end{align*}

\begin{enumerate}
    \item[(a)]
    \begin{proof}
        Let $P \curvearrowright N$ via conjugation, since $N \neq \{ e \}$ it is a p-group
        and thus the set of fixed points $N^P = 0 \mod p$ and since $e \in N^P$ this implies that
        $N^P \neq \{ e \}$. Now
        \[ N^P = \{ x \in N : \forall p \in P \text{ s.t. } pxp^{-1} = x \} = Z(P) \cap N \]
        hence $Z(P) \cap N \neq \{ e \}$.
    \end{proof}

    \item[(b)]
    \begin{proof}
        We proceed by induction. Namely we prove the following statement by induction on $j \in \mathbb N \cup \{ 0 \}$:
        If $P$ is a p-group such that $N \leq P$ and $p^j \mid |N|$ then $N$ has a normal subgroup of order $p^{j-1}$
        that is also normal in $P$.

        The base case is trivial as $p^0 = 1$ and the trivial subgroup $\{ e \}$ has the required properties.

        Now for the inductive step, suppose that the statement holds for some $j \in \mathbb N \cup \{ 0 \}$.
        Now consider a p-group $P$ and a normal subgroup $N$ such that $p^{j+1} \mid |N|$. Since $j \geq 0$
        this implies that $N \neq \{e \}$ and by $(a)$ this implies that $Z(P) \cap N$ is a non-trivial p-group.
        Thus by Cauchy's theorem there is $x \in Z(P) \cap N$ of order $p$. Now $\langle x \rangle \lhd N, P$.
        Moreover $P/\langle x \rangle$ is a p-group with normal subgroup $N/\langle x \rangle$ and $p^j \mid |N/\langle x \rangle|$.
        Thus by the inductive hypothesis $N/\langle x \rangle$ has a normal subgroup $\overline{H}$ of order $p^{j}$
        that is also normal in $P/\langle x \rangle$. By the correspondence isomorphism theorem,
        if $\pi$ is the natural projection for $P \to P/\langle x\rangle$ then $H = \pi^{-1}(\overline{H})$
        is a normal subgroup of $N$ of order $p^{j+1}$ which is also normal in $P$ as required.
    \end{proof}

    \item[(c)]
    \begin{proof}
        Suppose $G$ is a group, then for any prime $p$ applying $(b)$ to a Sylow p-subgroup
        of $G$ shows that for any $p^j \mid |G|$, $G$ has a subgroup of order $p^j$. However,
        such a subgroup need not be normal in $G$. For example take $S_3$ and the 2 element subgroup
        $\{ e, (12) \}$. This is not normal as
        \[ (23)(12)(23)^{-1} = (13) \notin \{ e, (12) \} \]
    \end{proof}

    \item[(d)]
    \begin{proof}
        Given a p-group $P$ with $|P| = p^a$, we prove the following stronger statement by induction on $i \in \mathbb N$:
        \[ |Z_i(P)| > p^{i-1} \]
        and when $i = a$ this implies that $|Z_a(P)| > p^{a-1}$ and hence $|Z_a(P)| = p^a$ by Lagrange's theorem
        which is equivalent to $Z_a(P) = P$ as required.

        For the base case $i = 1$: $|Z_1(P)| = |Z(P)| > p^0 = 1$ as a p-group has a non-trivial center.

        Now suppose the statement holds for some $i \in \{ 0, \ldots, a-1\}$. Then
        \[ Z_{i+1}(P) = \pi_i^{-1}(Z(P/Z_i(P))) \]
        hence
        \[ |Z_{i+1}(P)| = |Z_i(P)||Z(P/Z_i(P))| \]

        now first suppose that $Z(P/Z_i(P)) = \{ e \}$, this implies that $P/Z_i(P) = \{ e \}$
        as non-trivial p-groups have non-trivial centers. Thus $P = Z_i(P)$ and then clearly
        $Z_{i+1}(P) = P$ which clearly suffices.

        Otherwise, $Z(P/Z_i(P)) \neq \{ e \}$ and hence $|Z(P/Z_i(P))| \geq p$ and thus
        \[ |Z_{i+1}(P)| = |Z_i(P)||Z(P/Z_i(P))| > p^ip = p^{i+1} \]
        as required.
    \end{proof}

    \item[(e)]
    \begin{proof}
        To prove that a nilpotent group $G$ is solvable it suffices to show that
        \[ Z_{i+1}(G)/Z_{i}(G) \]
        is abelian then reversing the upper central series is clearly a solvable series for $G$.
        Now
        \[ Z_{i+1}(G) = \pi_i^{-1}(Z(G/Z_{i}(G))) \]
        so by definition
        \[ Z_{i+1}(G)/Z_i(G) = Z(G/Z_i(G)) \]
        hence $Z_{i+1}(G)/Z_i(G)$ is abelian and $G$ is solvable.

        For an example of a group that is solvable but not nilpotent consider $S_3$.

        $A_3 \cong C_3$, $S_3 \rhd A_3 \rhd \{e\}$ is a solvable series for $S_3$.
        However $Z(S_3) = \{e \}$ and thus $Z_1(S_3) = Z(S_3) = \{ e \}$ meaning
        that $Z_i(S_3) = \{e \}$ for all $i \in \mathbb N$. Hence $S_3$ is not nilpotent.
    \end{proof}
\end{enumerate}

\end{document}